% =============================================================================
%  macros.tex — пользовательские команды
% =============================================================================

% --- Пометки для работы над текстом ----------------------------------------
% Чтобы скрыть все пометки в финальной версии: \renewcommand{\TODO}[1]{}
\newcommand{\TODO}[1]{\textcolor{red}{\textbf{[TODO: #1]}}}
\newcommand{\NOTE}[1]{\textcolor{blue}{\textit{[Note: #1]}}}
\newcommand{\CITE}{\textcolor{orange}{\textbf{[cite]}}}
\newcommand{\NUM}{\textcolor{red}{\textbf{XX}}}   % заглушка для чисел

% --- Название системы (поменять в одном месте) ------------------------------
\newcommand{\sysname}{\textsc{DH-SLAM-1}\xspace}

% --- Сокращения ------------------------------------------------------------
\newcommand{\eg}{e.g.,\xspace}
\newcommand{\ie}{i.e.,\xspace}
\newcommand{\etal}{\textit{et al.}\xspace}
\newcommand{\wrt}{w.r.t.\xspace}

% --- Математические обозначения (SLAM) --------------------------------------
\newcommand{\R}{\mathbb{R}}
\newcommand{\SE}[1]{\mathrm{SE}(#1)}
\newcommand{\SO}[1]{\mathrm{SO}(#1)}
\newcommand{\pose}{\mathbf{T}}            % поза камеры
\newcommand{\rot}{\mathbf{R}}             % вращение
\newcommand{\trans}{\mathbf{t}}           % перенос
\newcommand{\landmark}{\mathbf{p}}        % 3D-точка карты
\newcommand{\meas}{\mathbf{z}}            % измерение
\newcommand{\state}{\mathbf{x}}           % вектор состояния
\newcommand{\jac}{\mathbf{J}}             % якобиан
\newcommand{\hess}{\mathbf{H}}            % (приближённый) гессиан
\newcommand{\infomat}{\boldsymbol{\Omega}}% информационная матрица
\newcommand{\norm}[1]{\left\lVert#1\right\rVert}
\DeclareMathOperator*{\argmin}{arg\,min}

% --- Аппаратные метрики -----------------------------------------------------
\newcommand{\fps}{\,FPS}
\newcommand{\mJpf}{\,\si{\milli\joule}/frame}
